\documentclass[11pt,letterpaper]{article}

\usepackage[spanish,es-noshorthands]{babel}
\usepackage{fontspec}
\setmainfont{Source Serif Pro}
\usepackage{microtype}
\usepackage{geometry}
\usepackage{booktabs}
\usepackage{tabularx}
\usepackage{enumitem}
\usepackage{xcolor}
\usepackage{graphicx}
\usepackage{csquotes}
\usepackage{amsmath}
\usepackage{siunitx}
\usepackage{tikz}
\usetikzlibrary{arrows.meta,positioning,shapes.geometric,fit,backgrounds,patterns}
\usepackage{xurl}
\usepackage[hidelinks]{hyperref}
\usepackage[backend=biber,style=apa,sorting=nyt]{biblatex}
\DeclareLanguageMapping{spanish}{spanish-apa}

\geometry{margin=2.3cm}
% Los rótulos de las figuras usan nodos de ancho fijo con texto alineado a la
% derecha, y la bibliografía compone en modo sloppy: ambos generan avisos de
% línea holgada que son cosméticos y no indican un problema de composición.
% Se silencian para que `make check` siga siendo sensible a los desbordes reales.
\hbadness=10000
\setlength{\emergencystretch}{1em}
\setlength{\parindent}{0pt}
\setlength{\parskip}{0.55em}
\setlist{nosep}
% Las figuras ocupan entre el 60 y el 70 por ciento de la caja de texto.
% Sin estos parámetros LaTeX las difiere todas al final del documento.
\renewcommand{\topfraction}{0.92}
\renewcommand{\bottomfraction}{0.85}
\renewcommand{\textfraction}{0.06}
\renewcommand{\floatpagefraction}{0.72}
\setcounter{topnumber}{2}
\setcounter{totalnumber}{3}
\addbibresource{referencias.bib}
\AtBeginBibliography{\sloppy}

\sisetup{
  output-decimal-marker = {,},
  group-separator = {.},
  group-minimum-digits = 4,
  detect-all
}

% Colores: sistema autónomo, conductor humano, regulador por homologación,
% regulador por autocertificación y señal de alerta
\definecolor{cADS}{RGB}{31,78,121}
\definecolor{cHum}{RGB}{191,110,30}
\definecolor{cReg}{RGB}{40,124,70}
\definecolor{cUSA}{RGB}{116,66,150}
\definecolor{cRiesgo}{RGB}{176,42,42}

% Barra horizontal con rótulo a la izquierda y valor a la derecha.
% #1 y, #2 fin de la barra, #3 color, #4 rótulo, #5 valor, #6 x del valor
\newcommand{\barrarot}[6]{%
  \node[left, text width=3.7cm, align=right, font=\scriptsize] at (-0.15,#1) {#4};%
  \fill[#3] (0,#1-0.17) rectangle (#2,#1+0.17);%
  \node[right, font=\scriptsize] at (#6,#1) {#5};%
}

\title{\vspace{-1.2cm}\textbf{Vehículos autónomos: certificar un proceso, no un
vehículo}\\
\large Cuatro preguntas con datos de agosto de 2026}
\author{Carlos Luis Rojas Aragonés\\
\small Gestión del Desarrollo Tecnológico}
\date{\small 19 de agosto de 2026}

\begin{document}

\maketitle
\vspace{-0.6em}

\section{Introducción}

Las cuatro preguntas del foro ya no son hipotéticas, y eso cambia cómo hay que
responderlas. A finales de marzo de 2026 el sistema de conducción de Waymo
acumulaba \textbf{\num{220.6} millones de millas recorridas sin conductor a
bordo} en cinco áreas metropolitanas de los Estados Unidos
\parencite{waymoimpact2026}, con del orden de medio millón de viajes pagados por
semana \parencite{cnbc2025rides}. El 31 de julio de 2026 la NHTSA concedió a Zoox
la primera exención de despliegue comercial para un vehículo construido en los
Estados Unidos sin volante ni pedales, con un límite de \num{2500} unidades al
año y dispensa de ocho normas federales \parencite{fr2026zoox}. El 4 de agosto de
2026 China publicó su norma obligatoria de seguridad para sistemas de conducción
automatizada, la GB 44721-2026, exigible desde el 1 de julio de 2027
\parencite{techtimes2026}. Y el Reino Unido tenía abierta, hasta el 9 de
septiembre de 2026, la consulta sobre el enunciado de principios de seguridad que
la Ley de Vehículos Automatizados de 2024 exige antes de autorizar el primer
vehículo \parencite{dft2026,ukavact2024}.

Mi respuesta a las cuatro preguntas parte de una sola idea, que anticipo aquí:
\textbf{lo que hay que certificar no es un vehículo, es un proceso}. La razón es
cuantitativa y la desarrollo en la sección siguiente: ninguna cantidad realista
de millas de prueba permite demostrar estadísticamente que un sistema de
conducción es seguro, y cada actualización de software reinicia esa cuenta. De
ahí se derivan las otras tres respuestas: la restricción que sirve no es la valla
física sino el dominio de diseño operativo autorizado; la responsabilidad tiene
que desplazarse del conductor a una entidad identificable que responda por cómo
conduce el vehículo; y el efecto sobre el empleo será pequeño en el agregado y
muy concentrado en el tiempo y en el territorio, lo que es precisamente lo que
hace que duela.

\section{Marco de análisis}

Uso tres conceptos y una unidad de medida.

\textbf{Nivel y dominio de diseño operativo.} La norma SAE J3016 define seis
niveles de automatización y, sobre todo, define el \emph{dominio de diseño
operativo} (ODD): el conjunto de condiciones (zona, vía, velocidad,
meteorología, hora) en las que el sistema está diseñado para funcionar
\parencite{sae2021}. Un nivel 4 no es \enquote{un coche que conduce solo}, es un
coche que conduce solo \emph{dentro de un ODD declarado}. Toda la discusión de
certificación y de zonas cerradas es, en realidad, una discusión sobre quién
autoriza y verifica ese ODD.

\textbf{Figura de mérito.} Una figura de mérito es una medida cuantitativa del
desempeño de una tecnología en la función que presta
\parencite{deweck2022}. Aquí la figura de mérito relevante no es la comodidad ni
el consumo, sino la tasa de siniestros con lesión por millón de millas, y su
comparación con la del conductor humano en la misma zona.

\textbf{El límite estadístico.} \textcite{kalra2016} calcularon cuántas millas
haría falta recorrer para demostrar que un sistema autónomo es más seguro que un
conductor humano. Para acreditar con un 95 por ciento de confianza una tasa de
1,09 muertes por cada 100 millones de millas hacen falta \textbf{275 millones de
millas sin una sola muerte}; para demostrar una mejora del 20 por ciento con el
mismo nivel de confianza, \textbf{\num{8800} millones de millas}. La conclusión
de los autores es que las pruebas de carretera, por sí solas, nunca serán
suficientes, y que la evaluación tiene que apoyarse en expedientes de seguridad,
simulación y auditoría del proceso de ingeniería, que es exactamente el enfoque
de la norma ISO 21448 sobre seguridad de la funcionalidad prevista y de la
ANSI/UL 4600 \parencite{iso21448,ul4600,koopman2017}.

\textbf{Y el marco jurídico.} En economía del derecho, la elección entre
responsabilidad por culpa y responsabilidad objetiva depende de quién controla el
riesgo y de quién puede reducirlo a menor coste \parencite{shavell1980}. Cuando
el conductor desaparece, el que controla el riesgo es quien escribe y actualiza
el software. La consecuencia es que el régimen tiene que virar de la culpa del
conductor a la responsabilidad del producto y del operador, y eso es literalmente
lo que están haciendo el Reino Unido y la Unión Europea.

\section{¿Cómo certificar y autorizar? El problema no cabe en millas}

\begin{figure}[tbp]
  \centering
  \begin{tikzpicture}
    % ---------------- panel A: umbrales de validación ----------------
    \node[font=\scriptsize\itshape, gray!35!black, right] at (0,4.10)
      {Panel A. Millas necesarias para demostrar la seguridad, en escala
       logarítmica};
    \foreach \e/\lab in {0/{$10^{8}$}, 1/{$10^{9}$}, 2/{$10^{10}$},
                         3/{$10^{11}$}, 4/{$10^{12}$}, 5/{$10^{13}$}} {
      \pgfmathsetmacro{\xg}{2.6*\e}
      \draw[gray!22] (\xg,-0.28) -- (\xg,3.30);
      \node[below, font=\tiny, gray!55!black] at (\xg,-0.28) {\lab};
    }
    \fill[cADS] (0,2.98) rectangle (0.894,3.32);
    \node[above right, font=\scriptsize, cADS] at (0,3.36)
      {Waymo sin conductor, acumulado a marzo de 2026};
    \node[right, font=\scriptsize] at (1.00,3.15) {\num{220.6} millones};
    \fill[cRiesgo] (0,1.98) rectangle (1.142,2.32);
    \node[above right, font=\scriptsize, cRiesgo] at (0,2.36)
      {umbral de Kalra y Paddock: acreditar la tasa de mortalidad al 95\,\%};
    \node[right, font=\scriptsize] at (1.25,2.15) {275 millones};
    \fill[cRiesgo!70] (0,0.98) rectangle (5.056,1.32);
    \node[above right, font=\scriptsize, cRiesgo] at (0,1.36)
      {umbral para demostrar una mejora del 20\,\% sobre el conductor humano};
    \node[right, font=\scriptsize] at (5.16,1.15) {\num{8800} millones};
    \fill[cHum] (0,-0.02) rectangle (11.758,0.32);
    \node[above right, font=\scriptsize, cHum] at (0,0.36)
      {millas que recorren los conductores humanos de EE. UU. en un año};
    \node[right, font=\scriptsize] at (11.86,0.15) {$3{,}3\times10^{12}$};

    % ---------------- panel B: trayectoria acumulada ----------------
    \begin{scope}[shift={(0,-6.70)}]
      \node[font=\scriptsize\itshape, gray!35!black, right] at (0,5.05)
        {Panel B. Millas acumuladas sin conductor y cruce del primer umbral};
      \foreach \yv/\lab in {0/0, 1.286/100, 2.571/200, 3.857/300} {
        \draw[gray!20] (0,\yv) -- (13.10,\yv);
        \node[left, font=\tiny, gray!55!black] at (-0.05,\yv) {\lab};
      }
      \node[rotate=90, font=\tiny, gray!55!black, align=center]
        at (-0.95,2.2) {millones de millas\\acumuladas};
      \draw[gray!60!black] (0,0) -- (13.10,0);
      \foreach \xv/\lab in {0/{2023}, 3.36/{2024}, 6.72/{2025}, 10.08/{2026}} {
        \draw[gray!60!black] (\xv,0) -- (\xv,-0.13);
        \node[below, font=\tiny, gray!55!black] at (\xv,-0.13) {\lab};
      }
      \draw[cRiesgo, thick, densely dotted] (0,3.536) -- (13.10,3.536);
      \node[above right, font=\tiny, cRiesgo] at (0.15,3.55)
        {275 millones: umbral de Kalra y Paddock};
      \draw[very thick, cADS] (3.08,0.091) -- (4.76,0.325) -- (6.44,0.643)
        -- (6.72,0.729) -- (8.40,1.286) -- (10.64,2.836);
      \foreach \x/\y in {3.08/0.091, 4.76/0.325, 6.44/0.643, 6.72/0.729,
                         8.40/1.286, 10.64/2.836} {
        \fill[cADS] (\x,\y) circle (1.8pt);
      }
      \node[right, font=\tiny, cADS] at (3.20,0.16) {7,1};
      \node[right, font=\tiny, cADS] at (4.88,0.22) {25,3};
      \node[left, font=\tiny, cADS] at (6.36,0.80) {50};
      \node[right, font=\tiny, cADS] at (6.86,0.66) {56,7};
      \node[right, font=\tiny, cADS] at (8.50,1.24) {100};
      \node[above left, font=\tiny, cADS] at (10.62,2.86) {220,6};
      \draw[thick, cADS, densely dashed] (10.64,2.836) -- (12.60,4.40);
      \draw[cRiesgo, thick] (11.52,3.536) circle (3.2pt);
      \draw[cRiesgo, thin] (11.60,3.40) -- (12.10,2.55);
      \node[font=\tiny, cRiesgo, align=center, below] at (12.10,2.55)
        {a 4 millones de millas\\[-0.3em] por semana, el umbral\\[-0.3em]
         se cruza hacia julio de 2026};
      \node[font=\tiny, cADS, rotate=38, above left] at (12.62,4.42)
        {extrapolación};
    \end{scope}
  \end{tikzpicture}
  \caption{El límite estadístico de la certificación por millas. Panel A: los
  dos umbrales de \textcite{kalra2016} frente a lo acumulado por el mayor
  operador del mundo y frente a la exposición anual de los conductores humanos
  de los Estados Unidos, calculada a partir de las \num{36640} muertes estimadas
  de 2025 y de la tasa de 1,10 por cada 100 millones de millas
  \parencite{nhtsa2026early}. Panel B: serie de millas acumuladas sin conductor
  a bordo, con los puntos documentados de 7,1 millones a finales de 2023
  \parencite{kusano2024}, 25,3 millones a mediados de 2024
  \parencite{dilillo2024}, 50 millones al cierre de 2024 y 100 millones en julio
  de 2025 \parencite{robotreport2025}, 56,7 millones del estudio revisado por
  pares \parencite{kusano2025} y \num{220.6} millones a marzo de 2026
  \parencite{waymoimpact2026}. La extrapolación usa el ritmo declarado de cuatro
  millones de millas semanales y es mía, no del operador. Elaboración propia.}
  \label{fig:validacion}
\end{figure}

La Figura~\ref{fig:validacion} contiene el argumento central de mi respuesta a
la primera pregunta. El operador más grande del mundo, después de una década de
desarrollo, está cruzando en estos meses el \emph{primero} de los dos umbrales
que hacen falta para una demostración estadística, y lo está cruzando en un ODD
muy concreto: cinco áreas metropolitanas de clima benigno. El segundo umbral, el
que permitiría afirmar con solidez estadística una mejora del 20 por ciento,
está cuarenta veces más lejos. Y hay un detalle que agrava el problema: cada
actualización relevante del software cambia el sistema evaluado, de modo que la
cuenta de millas no se acumula sobre un objeto fijo. Certificar por millas
recorridas es, por construcción, imposible.

Lo que sí se puede certificar es el proceso: el expediente de seguridad, la
trazabilidad de las decisiones de ingeniería, la validación en simulación, la
gestión de la seguridad de la organización y la vigilancia posterior al
despliegue. Es el enfoque de la ISO 21448 y de la UL 4600, y es también lo que ya
hace el Reglamento de Ejecución (UE) 2022/1426, que exige auditoría del sistema
de gestión de la seguridad y evaluación del concepto de seguridad del ADS ante la
autoridad de homologación, no una cifra de millas \parencite{eu2022reg}.

\begin{figure}[tbp]
  \centering
  \begin{tikzpicture}
    % ---------------- panel A: dos regímenes ----------------
    \node[right, text width=13.0cm, align=left, font=\scriptsize\itshape,
      gray!35!black] at (0,4.62)
      {Panel A. Dos maneras de autorizar, 2021 a 2027. Arriba, las
       jurisdicciones que homologan por tipo con auditoría previa; abajo, la
       autocertificación del fabricante con exenciones caso por caso};
    \draw[gray!60!black] (0,0.85) -- (13.20,0.85);
    \foreach \xv/\lab in {0/2021, 1.857/2022, 3.714/2023, 5.571/2024,
                          7.429/2025, 9.286/2026, 11.143/2027} {
      \draw[gray!60!black] (\xv,0.85) -- (\xv,0.72);
      \node[below, font=\tiny, gray!55!black, fill=white, inner sep=1.2pt]
        at (\xv,0.72) {\lab};
    }
    \fill[cReg!12] (11.143,0.92) rectangle (13.20,2.55);
    \node[font=\tiny, cReg, align=center] at (12.17,2.05)
      {2027: entran\\[-0.3em] en vigor el\\[-0.3em] enunciado\\[-0.3em]
       británico y la\\[-0.3em] norma china};
    % --- pista superior: homologación por tipo y auditoría ---
    \foreach \xv/\prof/\lab in {0.930/1.65/{Alemania: nivel 4\\con supervisor
                                  técnico},
                                3.170/0.55/{UE 2022/1426:\\homologación del ADS},
                                3.714/1.65/{R157 hasta\\130 km/h},
                                6.190/0.55/{Ley británica\\de 2024},
                                9.130/1.65/{primeros permisos\\de nivel 3 en
                                  China},
                                10.060/0.55/{consulta de\\principios (RU)},
                                10.370/2.70/{GB 44721-2026\\(China)}} {
      \draw[cReg, thick] (\xv,0.85) -- (\xv,0.85+\prof);
      \fill[cReg] (\xv,0.85+\prof) circle (1.6pt);
      \node[above, font=\tiny, cReg, align=center, fill=white,
        inner sep=1.5pt] at (\xv,0.92+\prof) {\lab};
    }
    \draw[cReg, thick] (0,0.85) -- (0,1.40);
    \fill[cReg] (0,1.40) circle (1.6pt);
    \node[above right, font=\tiny, cReg, align=left] at (-0.12,1.47)
      {UNECE R157:\\[-0.2em] mantenimiento\\[-0.2em] de carril};
    % --- pista inferior: autocertificación con exenciones ---
    \foreach \xv/\prof/\lab in {7.890/0.55/{marco de la NHTSA\\y AVEP ampliado},
                                9.590/1.55/{dos propuestas de\\reforma de las
                                  FMVSS},
                                10.210/0.55/{exención a Zoox\\y guía provisional}} {
      \draw[cUSA, thick] (\xv,0.85) -- (\xv,0.85-\prof);
      \fill[cUSA] (\xv,0.85-\prof) circle (1.6pt);
      \node[below, font=\tiny, cUSA, align=center, fill=white,
        inner sep=1.5pt] at (\xv,0.78-\prof) {\lab};
    }
    \draw[cReg, thick] (0.15,-1.85) -- (0.75,-1.85);
    \node[right, font=\tiny, cReg] at (0.85,-1.85)
      {homologación de tipo con auditoría};
    \draw[cUSA, thick] (5.30,-1.85) -- (5.90,-1.85);
    \node[right, font=\tiny, cUSA] at (6.00,-1.85)
      {autocertificación con exenciones};

    % ---------------- panel B: autorización por capas ----------------
    \begin{scope}[shift={(0,-7.60)}]
      \node[font=\scriptsize\itshape, gray!35!black, right] at (0,4.45)
        {Panel B. La autorización por capas que propondría};
      \foreach \i/\tit/\des in {%
        0/{1. Norma del vehículo}/{frenos, luces, visibilidad, protección de
           ocupantes: FMVSS o reglamentos UNECE, con exenciones motivadas y
           publicadas},
        1/{2. Expediente del ADS}/{concepto de seguridad auditado por un tercero
           acreditado: ISO 21448, UL 4600, anexo 3 del Reglamento (UE) 2022/1426},
        2/{3. ODD autorizado}/{zona, vías, velocidad, meteorología y horario
           concretos, con ampliación por tramos y revisión ante cada cambio de
           software},
        3/{4. Operación}/{supervisor técnico identificable, asistencia remota
           dimensionada, protocolo con bomberos y policía, formación del
           personal},
        4/{5. Vigilancia continua}/{telemetría y registro de datos obligatorios,
           publicación de siniestros, potestad de suspender la autorización sin
           tener que retirar el modelo}} {
        \pgfmathsetmacro{\yv}{0.80*\i}
        \fill[cReg!10] (0,\yv) rectangle (13.20,\yv+0.72);
        \fill[cReg] (0,\yv) rectangle (0.11,\yv+0.72);
        \node[right, text width=3.15cm, align=left, font=\scriptsize\bfseries,
          cReg!60!black] at (0.22,\yv+0.36) {\tit};
        \node[right, text width=8.55cm, font=\tiny, align=left]
          at (3.55,\yv+0.36) {\des};
      }
      \draw[-{Stealth[length=2mm]}, gray!60!black] (13.45,0.05) -- (13.45,4.25);
      \node[rotate=90, font=\tiny, gray!45!black, align=center]
        at (13.80,2.15) {del artefacto\\a la operación};
    \end{scope}
  \end{tikzpicture}
  \caption{Los dos regímenes vigentes y la propuesta. Panel A: en la pista
  superior, las jurisdicciones que autorizan por homologación de tipo con
  auditoría previa, con el Reglamento UNECE 157 de 2021 y su ampliación a
  \qty{130}{\kilo\meter\per\hour} de enero de 2023
  \parencite{unr157,unece2022}, el Reglamento de Ejecución (UE) 2022/1426
  \parencite{eu2022reg}, la ley alemana de conducción autónoma de 2021 con su
  reglamento de 2022, que exige un supervisor técnico
  \parencite{hoganlovells2021,taylorwessing2026}, la Ley británica de Vehículos
  Automatizados de 2024 con su consulta de principios de seguridad de junio de
  2026 \parencite{ukavact2024,dft2026} y los primeros permisos chinos de nivel 3
  de diciembre de 2025 seguidos de la norma obligatoria GB 44721-2026
  \parencite{cnevpost2025,techtimes2026}. En la pista inferior, el régimen
  estadounidense de autocertificación: el marco de la NHTSA de abril de 2025
  \parencite{nhtsa2025marco}, las dos propuestas de reforma de las FMVSS 102,
  103 y 104 de marzo de 2026 \parencite{sidley2026} y la exención a Zoox con la
  guía provisional del 31 de julio de 2026
  \parencite{fr2026zoox,fr2026guia}. Panel B: elaboración propia.}
  \label{fig:regimenes}
\end{figure}

Mi respuesta concreta a la primera pregunta es el Panel B de la
Figura~\ref{fig:regimenes}: \textbf{autorización por capas, con la última capa
siempre abierta}. Las cuatro primeras son las que ya existen, repartidas entre
regímenes distintos. La quinta es la que decide, y es la que menos se discute: si
la autoridad puede suspender una autorización de ODD en cuarenta y ocho horas
ante un patrón anómalo, sin tener que abrir un expediente de retirada del modelo,
entonces puede permitirse autorizar antes y con menos certeza. La reversibilidad
barata es lo que hace aceptable la incertidumbre. Ese es también el punto débil
del régimen estadounidense: funciona por exenciones puntuales de dos años
\parencite{fr2026zoox} y por guías no vinculantes \parencite{fr2026guia},
mientras la vigilancia depende de una orden general de comunicación de siniestros
cuyos informes, además, no determinan culpa \parencite{nhtsasgo2026}.

\section{¿Zonas cerradas o tránsito mixto? La valla útil ya existe y no es
física}

\begin{figure}[tbp]
  \centering
  \begin{tikzpicture}
    % ---------------- panel A: reducciones medidas ----------------
    \node[font=\scriptsize\itshape, gray!35!black, right] at (-3.9,3.95)
      {Panel A. Reducción de la tasa de choques frente al conductor humano de la
       misma zona, en \num{220.6} millones de millas};
    \foreach \xv/\lab in {0/0, 2/20, 4/40, 6/60, 8/80, 10/100} {
      \draw[gray!22] (\xv,-0.35) -- (\xv,3.45);
      \node[below, font=\tiny, gray!55!black] at (\xv,-0.35) {\lab\,\%};
    }
    \barrarot{3.10}{9.4}{cADS}{Lesión grave o mortal}{94\,\%}{9.5}
    \barrarot{2.48}{8.2}{cADS}{Despliegue de airbag}{82\,\%}{9.5}
    \barrarot{1.86}{8.2}{cADS}{Cualquier lesión}{82\,\%}{9.5}
    \barrarot{1.24}{9.3}{cADS!80}{Peatón con lesión}{93\,\%}{9.5}
    \barrarot{0.62}{8.4}{cADS!80}{Ciclista con lesión}{84\,\%}{9.5}
    \barrarot{0.00}{8.4}{cADS!80}{Motorista con lesión}{84\,\%}{9.5}

    % ---------------- panel B: quién provocó el choque ----------------
    \begin{scope}[shift={(0,-4.30)}]
      \node[font=\scriptsize\itshape, gray!35!black, right] at (-3.9,2.05)
        {Panel B. Los 78 siniestros con lesión o airbag comunicados entre agosto
         de 2025 y marzo de 2026};
      \node[left, text width=3.7cm, align=right, font=\scriptsize]
        at (-0.15,1.00) {Responsabilidad aparente};
      \fill[cHum] (0,0.83) rectangle (11.82,1.17);
      \fill[gray!45] (11.82,0.83) rectangle (12.47,1.17);
      \fill[cADS] (12.47,0.83) rectangle (12.80,1.17);
      \node[font=\tiny, white] at (5.9,1.00) {otro conductor: 72 de 78 (92\,\%)};
      \node[above, font=\tiny, gray!35!black] at (11.80,1.21) {4 dudosos};
      \node[above left, font=\tiny, cADS] at (12.80,1.48) {1 o 2 de Waymo};
      \draw[cADS, thin] (12.63,1.21) -- (12.68,1.46);
      \node[left, text width=3.7cm, align=right, font=\scriptsize]
        at (-0.15,0.15) {Tipo de impacto};
      \fill[cHum!85] (0,-0.02) rectangle (7.88,0.32);
      \fill[cHum!35] (7.88,-0.02) rectangle (12.80,0.32);
      \node[font=\tiny, white] at (3.94,0.15)
        {alcances por detrás: 48 de 78};
      \node[font=\tiny, gray!25!black] at (10.34,0.15)
        {resto: cruces, giros y adelantamientos, 30};
    \end{scope}
  \end{tikzpicture}
  \caption{La evidencia del tránsito mixto. Panel A: reducciones medidas por el
  operador sobre \num{220.6} millones de millas sin conductor, comparadas con la
  siniestralidad humana de las mismas zonas
  \parencite{waymoimpact2026,waymoblog2026}; el estudio revisado por pares con
  56,7 millones de millas obtiene reducciones del mismo orden desglosadas por
  tipo de choque \parencite{kusano2025}, y el análisis de reclamaciones de seguro
  de Swiss Re, del 88 por ciento en daños materiales y del 92 por ciento en daños
  personales sobre 25,3 millones de millas, es independiente del operador
  \parencite{dilillo2024}. Panel B: reparto de los 78 informes remitidos a la
  NHTSA, con la responsabilidad aparente leída de los relatos y sin resolución
  oficial de culpa, entre el 15 de agosto de 2025 y el 16 de marzo de 2026, según la lectura
  de los relatos hecha por \textcite{lee2026}; los informes de la orden general
  acreditan que hubo choque, no quién tuvo la culpa \parencite{nhtsasgo2026}.
  Elaboración propia.}
  \label{fig:mixto}
\end{figure}

Mi respuesta a la segunda pregunta es que \textbf{la separación física no es la
variable de control correcta, y además ya está superada por los hechos}. Tres
argumentos.

\textbf{Primero: el tránsito mixto ya se está midiendo, y el resultado no
sostiene la prohibición.} Sobre \num{220.6} millones de millas en tránsito
urbano abierto, con peatones, ciclistas y conductores humanos, las reducciones
son del 94 por ciento en choques con lesión grave o mortal y del 82 por ciento en
choques con cualquier lesión \parencite{waymoimpact2026}. Un estudio revisado por
pares con 56,7 millones de millas llega a magnitudes equivalentes desglosadas por
tipo de choque \parencite{kusano2025} y una aseguradora, con datos de
reclamaciones y no de la propia empresa, obtiene un 88 y un 92 por ciento
\parencite{dilillo2024}. Prohibir la mezcla exigiría demostrar que la mezcla es
peor que lo que ya tenemos, y la evidencia disponible apunta al contrario.

\textbf{Segundo: en la mezcla, el eslabón que falla es el humano.} De los 78
siniestros con lesión o airbag comunicados entre agosto de 2025 y marzo de 2026,
72 se originaron en la conducta de otro conductor y 48 fueron alcances por
detrás, muchos de ellos contra un vehículo detenido en un semáforo o cediendo el
paso a un peatón \parencite{lee2026}. Esto es coherente con el dato clásico de
que en el 94 por ciento de los choques la última causa de la cadena se atribuye
al conductor \parencite{singh2015}. Segregar los vehículos autónomos, en este
escenario, equivaldría a apartar al conductor que no comete la infracción.

\textbf{Tercero: la restricción sí existe, pero es el ODD.} Ningún robotaxi
circula hoy por donde quiere: circula por una zona geovallada, a velocidades
acotadas y con condiciones meteorológicas declaradas. Eso es una restricción real
y auditable, y tiene la propiedad de poder ampliarse por tramos, cosa que una
infraestructura dedicada no tiene: los carriles exclusivos exigen inversión
irreversible, no escalan al conjunto de la red y no resuelven el problema de las
intersecciones, que es donde ocurren los choques graves. Dicho esto, la
ampliación no puede ser automática. El caso del menor golpeado en un paso escolar
de Santa Mónica, investigado por la NHTSA desde enero de 2026
\parencite{cnbc2026}, es el tipo de suceso que justifica condiciones específicas
por microzona (entornos escolares, obras, eventos) antes que vallas generales.

\section{Implicaciones legales: la responsabilidad se traslada al que actualiza
el software}

\begin{figure}[tbp]
  \centering
  \begin{tikzpicture}[
      caja/.style={draw=gray!55, rounded corners=1.6pt, align=left,
        font=\tiny, text width=3.85cm, inner sep=3.2pt, minimum height=1.15cm},
      flecha/.style={-{Stealth[length=1.7mm]}, gray!55!black}]
    \foreach \i/\reg/\uno/\dos/\tres in {%
      0/{Dos sistemas autónomos entre sí}%
        /{Sin precedente publicado hasta agosto de 2026: los registros de las
          dos partes son la prueba}%
        /{Litigio o subrogación entre fabricantes y aseguradoras, sin conductor
          al que imputar}%
        /{El usuario queda fuera de la cadena; decide la telemetría, no el
          testimonio},
      1/{Unión Europea y Alemania}%
        /{Seguro obligatorio del titular: indemniza a la víctima}%
        /{Acción contra el fabricante por producto defectuoso, Directiva
          2024/2853}%
        /{Presunción de defecto y exhibición de pruebas si la complejidad
          técnica lo exige},
      2/{Reino Unido}%
        /{La aseguradora del vehículo paga en primera instancia, ley de 2018}%
        /{Repetición contra la entidad autorizada de conducción, ley de 2024}%
        /{El usuario al mando no responde de cómo conduce el vehículo},
      3/{Estados Unidos}%
        /{Demanda civil de la víctima: negligencia y producto defectuoso}%
        /{El jurado reparte porcentajes de culpa entre las partes}%
        /{\emph{Benavides contra Tesla} (2025): 33\,\% al fabricante, en
          apelación}} {
      \pgfmathsetmacro{\yv}{2.05*\i}
      \node[right, font=\scriptsize\bfseries, cADS] at (0,\yv+1.55) {\reg};
      \node[caja, anchor=south west] at (0,\yv) {\uno};
      \node[caja, anchor=south west] at (4.55,\yv) {\dos};
      \node[caja, anchor=south west] at (9.10,\yv) {\tres};
      \draw[flecha] (4.00,\yv+0.575) -- (4.50,\yv+0.575);
      \draw[flecha] (8.55,\yv+0.575) -- (9.05,\yv+0.575);
    }
    \node[font=\tiny, gray!45!black, above right] at (0,8.30)
      {quién indemniza primero $\rightarrow$ contra quién se repite
       $\rightarrow$ qué decide el caso};
  \end{tikzpicture}
  \caption{Quién responde en cada régimen. Estados Unidos no tiene una regla
  especial para vehículos autónomos: la asignación la hace el jurado caso por
  caso, como en el veredicto de \emph{Benavides contra Tesla} de agosto de 2025,
  que atribuyó al fabricante el 33 por ciento de la responsabilidad y está
  recurrido \parencite{wshb2025}. El Reino Unido separa el pago de la
  imputación: la aseguradora indemniza de entrada \parencite{ukaeva2018} y luego
  repite contra la entidad autorizada de conducción, que es la responsable de
  cómo conduce el vehículo \parencite{ukavact2024,cms2024}. En la Unión Europea
  el seguro obligatorio del titular convive con la nueva directiva de productos
  defectuosos, que incluye el software y las actualizaciones, invierte de hecho
  la carga de la prueba cuando la complejidad técnica lo justifica y obliga a
  exhibir pruebas \parencite{eupld2024,reedsmith2025}; Alemania añade la figura
  del supervisor técnico, cubierta por el seguro
  \parencite{hoganlovells2021,taylorwessing2026}. Elaboración propia.}
  \label{fig:responsabilidad}
\end{figure}

La tercera pregunta tiene, a mi juicio, una respuesta clara en cuanto al
principio y complicada en cuanto a la prueba.

\textbf{El principio: donde no hay conductor, no puede haber culpa del
conductor.} Quien controla el riesgo es quien diseña, entrena y actualiza el
sistema, de modo que la responsabilidad debe recaer sobre el operador y el
fabricante, no sobre el ocupante \parencite{shavell1980}. El Reino Unido lo ha
resuelto de la forma más limpia: el usuario al mando no responde de cómo conduce
el vehículo, la aseguradora indemniza en primera instancia y luego repite contra
la entidad autorizada de conducción \parencite{ukaeva2018,ukavact2024,cms2024}. La
víctima no tiene que averiguar quién falló para cobrar, que es exactamente lo que
hay que conseguir.

\textbf{Contra un conductor convencional} el resultado depende del régimen. En la
Unión Europea el seguro obligatorio del titular paga y la acción posterior se
dirige contra el fabricante por producto defectuoso, con dos novedades que
importan mucho: la Directiva (UE) 2024/2853 considera producto al software y a
sus actualizaciones, y permite presumir el defecto y el nexo causal cuando la
complejidad técnica hace excesivamente difícil la prueba, además de obligar al
demandado a exhibir la información técnica \parencite{eupld2024,reedsmith2025}.
En los Estados Unidos, sin regla especial, el reparto lo hace un jurado, con la
volatilidad que muestra el 33 por ciento imputado al fabricante en
\emph{Benavides contra Tesla} \parencite{wshb2025}.

\textbf{Entre dos vehículos autónomos} no hay precedente publicado hasta agosto
de 2026, y no es casualidad: son pocos, están concentrados y sus siniestros con
otros vehículos son casi siempre alcances provocados por humanos
\parencite{lee2026,dmv2026}. Cuando ocurra, el caso se decidirá por los registros
de datos de los dos sistemas, no por testimonios, y se resolverá entre
fabricantes y aseguradoras sin que ningún usuario tenga nada que declarar. Esto
tiene una consecuencia práctica que sí depende de nosotros: \textbf{el registro
de datos del suceso tiene que ser obligatorio, normalizado e independiente del
fabricante}. Sin eso, el régimen de responsabilidad es teórico. Los informes de
la orden general estadounidense sirven para vigilar tendencias, pero
explícitamente no determinan culpa \parencite{nhtsasgo2026}, y el propio Cuadro
de siniestros del regulador californiano no resuelve responsabilidades
\parencite{dmv2026}.

\section{¿Empleo? Poco en el agregado, mucho en cada barrio afectado}

\begin{figure}[tbp]
  \centering
  \begin{tikzpicture}
    % ---------------- panel A: empleo expuesto ----------------
    \node[font=\scriptsize\itshape, gray!35!black, right] at (-3.9,2.70)
      {Panel A. Empleo expuesto en los Estados Unidos, en millones de personas};
    \foreach \xv/\lab in {0/0, 2.4/4, 4.8/8, 7.2/12, 9.6/16} {
      \draw[gray!22] (\xv,-0.30) -- (\xv,2.30);
      \node[below, font=\tiny, gray!55!black] at (\xv,-0.30) {\lab};
    }
    \barrarot{1.80}{9.30}{cHum!45}{Ocupaciones afectadas en algún
      grado}{15,5}{9.45}
    \barrarot{1.00}{2.28}{cHum!75}{Operadores de vehículos de
      motor}{3,8}{2.42}
    \barrarot{0.20}{1.34}{cHum}{Camión de gran tonelaje y
      tractocamión}{2,24}{1.48}
    \node[right, text width=9.6cm, font=\tiny, align=left, gray!25!black]
      at (-0.10,-0.98) {y aun así la proyección oficial para 2024 a 2034 es de
      crecimiento: +4\,\% en camión de gran tonelaje, +8\,\% en reparto y
      +9\,\% en taxi y vehículo con conductor};

    % ---------------- panel B: efectos ya medidos ----------------
    \begin{scope}[shift={(0,-5.05)}]
      \node[font=\scriptsize\itshape, gray!35!black, right] at (-3.9,2.90)
        {Panel B. Variación interanual en el cuarto trimestre de 2025, en puntos
         porcentuales};
      \foreach \xv/\lab in {0/{$-6$}, 1.2/{$-4$}, 2.4/{$-2$}, 3.6/{0},
                            4.8/{$+2$}} {
        \draw[gray!22] (\xv,-0.55) -- (\xv,2.50);
        \node[below, font=\tiny, gray!55!black] at (\xv,-0.55) {\lab};
      }
      \draw[gray!60!black, thick] (3.6,-0.55) -- (3.6,2.50);
      \foreach \y/\x/\lab/\v/\col in {%
        2.10/4.68/{Ingreso por hora,\\media nacional}/{$+1{,}8$}/{cReg},
        1.35/3.36/{Ingreso por hora,\\Phoenix}/{$-0{,}4$}/{cRiesgo!70},
        0.60/2.58/{Ingreso por hora,\\Los Ángeles}/{$-3{,}7$}/{cRiesgo},
        -0.15/1.62/{Viajes por hora,\\mercados con robotaxi}/{$-5{,}3$}/{cRiesgo}} {
        \node[left, text width=3.7cm, align=right, font=\scriptsize]
          at (-0.15,\y) {\lab};
        \fill[\col] (3.6,\y-0.16) rectangle (\x,\y+0.16);
        \node[right, font=\scriptsize] at (4.95,\y) {\v};
      }
      \node[right, text width=5.9cm, font=\tiny, align=left, gray!25!black]
        at (5.85,1.60) {para comparar, los viajes por hora cayeron un
        2,6\,\% en el conjunto del país: la mitad de la caída registrada en
        las ciudades con robotaxi};
    \end{scope}
  \end{tikzpicture}
  \caption{Empleo expuesto y efectos ya observables. Panel A: \num{15.5}
  millones de trabajadores en ocupaciones afectadas en algún grado y 3,8
  millones de operadores de vehículos de motor \parencite{beede2017}, frente a
  los \num{2235100} conductores de camión de gran tonelaje contabilizados en 2024
  y a proyecciones oficiales que siguen siendo de crecimiento
  \parencite{bls2026camion,bls2026taxi}. Panel B: variación interanual del cuarto
  trimestre de 2025 en los mercados con robotaxi frente al conjunto del país,
  según los datos propios de una plataforma de análisis para conductores
  \parencite{gridwise2026}; es un panel no representativo y así hay que leerlo.
  Elaboración propia.}
  \label{fig:empleo}
\end{figure}

Creo que la pregunta está mal planteada cuando se formula como pérdida
\emph{o} incremento, porque las dos cosas ocurren y no a la misma velocidad.

\textbf{En el agregado, el efecto medido es pequeño.} El estudio más cuidadoso
que conozco estima entre 1,3 y 2,3 millones de empleos afectados en treinta años
en los Estados Unidos, con un efecto sobre la tasa de desempleo de entre 0,06 y
0,13 puntos porcentuales en la década de mayor impacto \parencite{groshen2018}.
Esa magnitud es la mitad de lo que supuso la pérdida de empleo industrial y está
repartida en tres décadas. La proyección oficial de empleo, entretanto, sigue
siendo de crecimiento en las tres ocupaciones de conducción más numerosas
\parencite{bls2026camion,bls2026taxi}, lo que da una idea de lo lento que el
regulador espera que sea el proceso.

\textbf{En lo local, el efecto ya se nota.} En el cuarto trimestre de 2025, los
viajes por hora de los conductores de plataforma cayeron un 5,3 por ciento en los
mercados con robotaxi frente a un 2,6 por ciento en el conjunto del país, y el
ingreso por hora bajó en Los Ángeles, San Francisco y Phoenix mientras subía un
1,8 por ciento a escala nacional \parencite{gridwise2026}. Es exactamente el
patrón que describe la literatura de automatización: el efecto de
desplazamiento llega antes y más concentrado que el de recomposición
\parencite{acemoglu2020,autor2015}.

\textbf{Y los empleos nuevos existen, pero no son los mismos ni están en el
mismo sitio.} La operación de una flota sin conductor necesita asistencia remota,
mantenimiento, limpieza, gestión de depósitos, cartografía, anotación de datos y
respuesta a incidentes. Cruise declaró un asistente remoto por cada 15 a 20
vehículos \parencite{cnbc2023cruise} y Alemania exige por ley una supervisión
técnica humana capaz de detener el vehículo y hablar con los ocupantes
\parencite{hoganlovells2021,taylorwessing2026}. Son puestos reales, pero con una
razón de sustitución desfavorable, con otro perfil de cualificación y
concentrados en las sedes del operador, no en las ciudades donde se pierde el
turno de taxi.

\section{Cronología y matriz de respuestas}

El Cuadro~\ref{tab:crono} ordena los hechos que sostienen las cuatro
respuestas y el Cuadro~\ref{tab:matriz} las resume en forma de matriz, por
escenario de siniestro y por régimen jurídico.

\begin{table}[tbp]
  \centering
  \small
  \caption{Cronología verificada del proceso regulatorio y del despliegue.}
  \label{tab:crono}
  \begin{tabularx}{\textwidth}{@{}l >{\raggedright\arraybackslash}X@{}}
    \toprule
    \textbf{Fecha} & \textbf{Hecho} \\
    \midrule
    01/2021 & Entra en vigor el Reglamento UNECE 157 para el mantenimiento
      automatizado de carril \parencite{unr157} \\
    04/2021 & SAE J3016: última revisión de la taxonomía de niveles y del
      dominio de diseño operativo \parencite{sae2021} \\
    07/2021 & Alemania aprueba la ley de conducción autónoma de nivel 4 con
      supervisor técnico \parencite{hoganlovells2021} \\
    09/2022 & Reglamento de Ejecución (UE) 2022/1426: homologación del ADS con
      auditoría del sistema de gestión \parencite{eu2022reg} \\
    01/2023 & El Reglamento 157 se amplía a \qty{130}{\kilo\meter\per\hour} y
      permite cambio de carril \parencite{unece2022} \\
    05/2024 & Ley británica de Vehículos Automatizados: entidad autorizada de
      conducción y usuario al mando \parencite{ukavact2024} \\
    10/2024 & Directiva (UE) 2024/2853: el software es producto y se presume el
      defecto si la prueba es excesivamente difícil \parencite{eupld2024} \\
    24/04/2025 & La NHTSA publica su marco de vehículos autónomos y amplía el
      programa de exenciones a los vehículos nacionales
      \parencite{nhtsa2025marco} \\
    15/07/2025 & Waymo alcanza 100 millones de millas sin conductor, el doble que
      seis meses antes \parencite{robotreport2025} \\
    01/08/2025 & Veredicto en \emph{Benavides contra Tesla}: 33 por ciento de
      responsabilidad al fabricante, recurrido \parencite{wshb2025} \\
    12/2025 & China concede los primeros permisos de nivel 3 a turismos, en
      rutas designadas \parencite{cnevpost2025} \\
    08/12/2025 & Waymo supera los \num{450000} viajes pagados por semana
      \parencite{cnbc2025rides} \\
    29/01/2026 & La NHTSA abre investigación por el atropello de un menor en un
      entorno escolar de Santa Mónica \parencite{cnbc2026} \\
    16/03/2026 & Dos propuestas de reforma de las FMVSS 102, 103 y 104 para
      vehículos sin mandos manuales \parencite{sidley2026} \\
    31/03/2026 & \textbf{\num{220.6} millones de millas} acumuladas sin conductor
      a bordo \parencite{waymoimpact2026} \\
    17/06/2026 & El Reino Unido abre la consulta del enunciado de principios de
      seguridad, con plazo hasta el 9 de septiembre \parencite{dft2026} \\
    31/07/2026 & Primera exención de despliegue comercial a un vehículo sin
      volante fabricado en EE. UU., hasta \num{2500} unidades al año
      \parencite{fr2026zoox} \\
    04/08/2026 & China publica la norma obligatoria GB 44721-2026, exigible desde
      julio de 2027 \parencite{techtimes2026} \\
    2.º sem. 2027 & Fecha prevista de plena aplicación del régimen británico
      \parencite{dft2026} \\
    \bottomrule
  \end{tabularx}
\end{table}

\begin{table}[tbp]
  \centering
  \small
  \caption{Quién responde según el escenario y el régimen.}
  \label{tab:matriz}
  \begin{tabularx}{\textwidth}{@{}>{\raggedright\arraybackslash}p{0.19\textwidth}
    >{\raggedright\arraybackslash}X >{\raggedright\arraybackslash}X
    >{\raggedright\arraybackslash}X@{}}
    \toprule
    \textbf{Escenario} & \textbf{Estados Unidos} & \textbf{Reino Unido}
      & \textbf{UE y Alemania} \\
    \midrule
    Autónomo contra conductor humano
      & jurado: negligencia y producto, reparto por porcentajes
      & aseguradora primero, repetición contra la entidad autorizada
      & seguro del titular primero, luego producto defectuoso \\
    Autónomo contra autónomo
      & sin precedente: litigio entre fabricantes
      & dos entidades autorizadas frente a frente
      & dos fabricantes, con exhibición de pruebas obligatoria \\
    Atropello de peatón
      & responsabilidad civil, con riesgo de daños punitivos
      & igual que el primer escenario, sin culpa del usuario
      & igual, más responsabilidad penal del supervisor si desatiende \\
    Nivel 3 con conductor presente
      & zona gris: la transferencia de control decide el caso
      & el usuario al mando responde solo fuera del modo autónomo
      & el conductor responde tras la petición de recuperación \\
    Fallo de la asistencia remota
      & operador, por negligencia en la operación
      & la entidad autorizada, por incumplir sus obligaciones
      & el supervisor técnico y el operador, con seguro obligatorio \\
    Prueba que decide
      & la que consiga el litigio
      & los registros exigidos a la entidad autorizada
      & registro del suceso más presunción a favor de la víctima \\
    \bottomrule
  \end{tabularx}
\end{table}

\section{Precisiones y limitaciones}

\begin{enumerate}
  \item \textbf{Las reducciones del 94 y del 82 por ciento las publica el propio
    operador.} Están construidas con una referencia humana ajustada a la misma
    zona y tipo de vía, hay una versión revisada por pares con 56,7 millones de
    millas \parencite{kusano2025} y hay una tercera fuente independiente, la
    aseguradora \parencite{dilillo2024}, pero no hay una auditoría pública de los
    denominadores. Lo trato como evidencia fuerte, no como dato oficial.
  \item \textbf{El ODD no es comparable con el del conductor humano.} Las millas
    acumuladas se han recorrido en cinco áreas de clima mayoritariamente benigno
    y a velocidades urbanas. Nada de esto acredita desempeño en nieve, en
    carretera secundaria sin señalizar ni en autopista a
    \qty{120}{\kilo\meter\per\hour}.
  \item \textbf{El reparto de culpa del Panel B de la Figura~\ref{fig:mixto} es
    una lectura de relatos, no una resolución.} La hizo un analista externo sobre
    los informes remitidos al regulador \parencite{lee2026}, y esos informes
    acreditan que hubo choque, no quién tuvo la culpa \parencite{nhtsasgo2026}.
  \item \textbf{Los umbrales de Kalra y Paddock suponen un sistema fijo y
    ausencia total de fallos.} Ninguna de las dos condiciones se cumple: hay
    siniestros y hay actualizaciones frecuentes. Por eso los uso como argumento
    de imposibilidad de la certificación por millas, no como una meta que
    alguien deba alcanzar.
  \item \textbf{Los datos de empleo de plataforma proceden de un panel privado.}
    Los porcentajes de \textcite{gridwise2026} vienen de los conductores que usan
    una aplicación concreta, con sesgo de selección desconocido y sin ajuste por
    la evolución de la demanda. Marcan una dirección, no una magnitud.
  \item \textbf{Los dos fallecimientos con vehículos sin conductor implicados no
    se atribuyeron al sistema.} En el choque múltiple de San Francisco de enero
    de 2025 y en el atropello de Dallas de 2026, la reconstrucción sitúa la causa
    en terceros \parencite{nbcbay2025,templeton2026}. Lo señalo porque la
    ausencia de una muerte imputada al sistema es la razón por la que el debate
    sigue siendo técnico y no político, y eso puede cambiar con un solo caso.
  \item \textbf{No hay contraste posible del escenario de dos autónomos.} Mi
    respuesta a esa parte de la pregunta es un razonamiento sobre los regímenes
    vigentes, no una descripción de la práctica: la práctica todavía no existe.
  \item \textbf{China entra en el análisis con fuentes de prensa
    especializada.} Ni el texto de la GB 44721-2026 ni los permisos de nivel 3
    están disponibles en una versión oficial que yo pueda verificar, de modo que
    esos dos hitos tienen menos solidez que los europeos y estadounidenses
    \parencite{cnevpost2025,techtimes2026}.
\end{enumerate}

\section*{Conclusiones}

\textbf{Certificación.} Por proceso y por capas, no por millas ni por modelo. El
umbral estadístico que haría creíble una certificación basada en kilometraje está
fuera de alcance \parencite{kalra2016}, de modo que hay que auditar el expediente
de seguridad, autorizar un ODD explícito, exigir supervisión y telemetría y
conservar la potestad de suspender esa autorización de un día para otro. La
homologación europea con auditoría se acerca más a ese modelo que la
autocertificación con exenciones puntuales.

\textbf{Zonas cerradas o tránsito mixto.} Mixto, con el ODD como valla. La
segregación física es caro, irreversible y no ataca el punto donde se producen
los choques graves; la geovalla, en cambio, es una restricción real que se puede
ampliar por tramos y revocar. La evidencia disponible sobre \num{220.6} millones
de millas en tránsito abierto no sostiene la prohibición de mezclar, y sí sostiene
condiciones específicas para microzonas sensibles.

\textbf{Responsabilidad.} Del conductor al operador y al fabricante, con pago en
primera instancia por el seguro y repetición posterior, que es el diseño
británico. Entre dos vehículos autónomos, la disputa será entre fabricantes y se
decidirá por los registros de datos, lo que convierte al registro obligatorio y
normalizado en el requisito jurídico más importante de todo el paquete.

\textbf{Empleo.} Las dos cosas, pero desalineadas. Poco en el agregado y en
treinta años \parencite{groshen2018}, mucho en unas pocas ciudades y en dos o
tres años, lo que ya se está midiendo \parencite{gridwise2026}. Los puestos
nuevos son reales y necesarios, pero con razones de uno por cada quince o veinte
vehículos y con otro perfil.

La lección que me llevo como CTO es de gobierno de la tecnología, no de
tecnología: \textbf{cuando no se puede demostrar que un sistema es seguro antes
de desplegarlo, la única alternativa responsable es hacer que desplegarlo sea
reversible}. Autorización acotada, medición obligatoria, publicación de los
siniestros y capacidad de apagarlo. Es menos elegante que una certificación
definitiva, y es lo único que funciona cuando el sistema que hay que certificar
cambia con cada actualización.

\printbibliography[title={Referencias bibliográficas}]

\end{document}
